\chapter{Operands}
\label{ch:operands}

An instruction's operands are separated by commas, with the destination first. This chapter lists every operand form
the assembler accepts. The entry for each instruction in the \archmanual{} gives the forms that instruction accepts,
and the Addressing Modes chapter of the \archmanual{} gives the meaning of each addressing mode.

\section{Operand Forms}

\begin{table}[H]
	\centering
	\begin{tabular}{lll}
		\toprule
		\textbf{Form}                                  & \textbf{Name}                     & \textbf{Example}              \\
		\midrule
		\texttt{\#n}                                   & Immediate number                  & \texttt{ADDI r1, \#12}        \\
		\texttt{@label}                                & Immediate label reference         & \texttt{MOVI al, @table}      \\
		\texttt{\$NAME}                                & Immediate named constant          & \texttt{CMPI r1, \$LIMIT}     \\
		\texttt{r1}--\texttt{r7}, \texttt{sr}, halves  & Register direct                   & \texttt{MOVR r2, r1}          \\
		\texttt{a}, \texttt{s}, \texttt{p}, \texttt{l} & Address register pair             & \texttt{LJMP a}               \\
		\texttt{(addr)}                                & Indirect, zero displacement       & \texttt{LOAD r1, (a)}         \\
		\texttt{(\#n, addr)}                           & Indirect, immediate offset        & \texttt{LOAD r1, (\#4, a)}    \\
		\texttt{(rO, addr)}                            & Indirect, register offset         & \texttt{STOR (r2, a), r1}     \\
		\texttt{(addr)+}                               & Post-increment, zero displacement & \texttt{LOAD r1, (a)+}        \\
		\texttt{(\#n, addr)+}                          & Post-increment, immediate         & \texttt{LOAD r1, (\#0, a)+}   \\
		\texttt{(rO, addr)+}                           & Post-increment, register          & \texttt{LOAD r1, (r2, s)+}    \\
		\texttt{-(addr)}                               & Pre-decrement, zero displacement  & \texttt{STOR -(s), r1}        \\
		\texttt{-(\#n, addr)}                          & Pre-decrement, immediate          & \texttt{STOR -(\#0, s), r1}   \\
		\texttt{-(rO, addr)}                           & Pre-decrement, register           & \texttt{STOR -(r2, s), r1}    \\
		\texttt{type \#count}, \texttt{type rC}        & Shift                             & \texttt{ADDR r1, r2, ASL \#1} \\
		\bottomrule
	\end{tabular}
	\caption{Operand Forms}
	\label{tab:operand-forms}
\end{table}

In the indirect forms \texttt{addr} is one of the four address register pairs, \texttt{a}, \texttt{s}, \texttt{p} or
\texttt{l}. The displacement is added to the low word of the pair. The offset in parentheses may be a number, a label
reference or a named constant.

An immediate displacement of zero may be omitted: \texttt{(addr)}, \texttt{(addr)+} and \texttt{-(addr)} are the same
as \texttt{(\#0, addr)}, \texttt{(\#0, addr)+} and \texttt{-(\#0, addr)}. The assembler accepts the shorthand only
where it accepts the full form, and \mnemonic{LJMP} and \mnemonic{LJSR} take \texttt{addr} rather than \texttt{(addr)}.
Register displacements have no shorthand.

\section{Immediate Operands}

An immediate operand may be a number, a label reference, or a named constant. A number is written with a leading hash
(Table~\ref{tab:number-forms}); a label reference and a named constant are not.

A label reference is accepted only by \mnemonic{MOVI}, \mnemonic{LOAD}, \mnemonic{STOR}, \mnemonic{LDEA},
\mnemonic{LDEL}, \mnemonic{LJMP}, \mnemonic{LJSR}, \mnemonic{BRAN} and \mnemonic{BRSR}. Other instructions with an
immediate operand accept a number or a named constant only. To use an address with one of them, load it with
\mnemonic{MOVI} first.

An immediate operand has the width of the instruction format it is encoded in: 16 bits in the immediate format and 8
bits in the short immediate format. The assembler is not required to diagnose a value that is too wide for its field,
so a value must be kept in range.

\section{Label References}

A label reference in an operand is resolved by the linker, not the assembler. What it resolves to depends on the
instruction and on an optional suffix; Chapter~\ref{ch:directives-and-symbols} gives the rules.

\section{Shift Operands}

A shift is written as a shift type followed by a count. The type is one of \texttt{NUL}, \texttt{LSL}, \texttt{LSR},
\texttt{ASL}, \texttt{ASR}, \texttt{RTL} or \texttt{RTR}. The count is either an immediate number from 0 to 15 or a
register that holds the count. A space separates the type from the count.

% qa:assemble
\begin{lstlisting}
    ADDR r1, r2, LSL #2      ; literal count
    ADDR r1, r2, LSR r3      ; count held in r3
\end{lstlisting}

A shift applies to one operand only, as defined for each instruction in the \archmanual{}. A count above 15 is rejected
by the assembler. The Shift Operations chapter of the \archmanual{} defines each shift type.

\section{Register Shorthand}

A save-result register-form ALU instruction, such as \mnemonic{ADDR}, written with two register operands is shorthand
for the three-operand form with the destination repeated as the first source. \mnemonic{MOVR} is the exception: its
only form is \texttt{MOVR rD, rS}. These two lines assemble to the same instruction:

% qa:assemble
\begin{lstlisting}
    ADDR r1, r1, r2         ; r1 = r1 + r2
    ADDR r1, r2             ; r1 = r1 + r2
\end{lstlisting}

A shift written with the shorthand applies to the left operand, which is the destination register, not to the second
register written.

The test instructions \mnemonic{CMNR}, \mnemonic{CNCR}, \mnemonic{CMPR}, \mnemonic{CPCR}, \mnemonic{TSAR},
\mnemonic{TSOR} and \mnemonic{TSXR} have no destination. Each takes exactly two registers, \texttt{rS1, rS2\{,
	shift\}}.
